\title{Byte Latent Transformer: Patches Scale Better Than Tokens}

\begin{abstract}
We present the Byte Latent Transformer ({\textsc BLT}), an innovative byte-level LLM architecture that, for the first time, achieves parity with tokenization-based LLMs at scale while delivering considerable gains in inference speed and resilience. {\textsc BLT} represents byte sequences as variable-length patches, treating these patches as its core computational units. The model dynamically segments these patches according to the predicted entropy of the upcoming byte, allocating greater computational effort and model resources to segments characterized by heightened data complexity. We conduct the inaugural scaling analysis of byte-level models regulated by FLOPs, evaluating configurations up to 8B parameters trained on 4T bytes. Our empirical findings validate the practicality of scaling byte-level models without relying on a predetermined vocabulary. Dynamic patch length selection, tailored to data predictability, leads to enhanced efficiency in both training and inference, and delivers observed qualitative gains in reasoning and rare-case generalization. Ultimately, at a constant inference budget, {\textsc BLT} achieves superior scaling relative to models that depend on tokenization, by jointly increasing both the patch length and the overall model capacity.
\end{abstract}

\section{Introduction}
\label{section:intro}

We present the Byte Latent Transformer~(\textbf{BLT}), an architecture that eliminates the need for tokenization by processing raw byte sequences directly. This approach is the first to demonstrate performance on par with conventional models that rely on tokenization, even at large scales, while also offering notable gains in both efficiency and robustness (\S\ref{sec:robustness}). While most large language models (llms) are built to be almost entirely end-to-end, tokenization remains a handcrafted pre-processing step that clusters byte sequences into a predefined set of tokens. This static mapping introduces biases in string compression, resulting in various limitations—such as heightened sensitivity to the choice of domain or modality~\citep{dagangetting}, increased vulnerability to noisy input (\S\ref{sec:robustness}), insufficient modeling of orthography~\citep{edman2024cute}, and disparities in multilingual processing~\citep{liang2023xlm,petrov2024language,limisiewicz2024myte}.

Historically, tokenization has played a crucial role because training llms directly on byte-level inputs is highly resource-intensive at large scales, mainly due to the extended sequence lengths involved~\citep{xue2022byt5}. Some earlier studies have sought to alleviate this challenge by introducing more efficient forms of self-attention~\citep{el2020characterbert, clark2022canine} or by leveraging attention-free model designs~\citep{wang2024mambabyte}~(\S\ref{section:related_work}). Nonetheless, these strategies mostly benefit the training of \textit{small models}. When scaling up, the bulk of computational overhead in Transformers arises from processing vast feed-forward network layers across every byte, rather than from the self-attention computation.

For more effective compute distribution, we introduce an adaptive, learnable approach that organizes bytes into \textit{patches}~(\S\ref{section:patching}) as well as a novel model architecture that integrates information at both the byte and patch levels. In contrast to tokenization, {\textsc BLT} does not employ a predetermined vocabulary for patches. Instead, any combination of bytes can be transformed into latent patch representations using efficient, parameterized encoder and decoder components. Our experiments demonstrate that this methodology enables a \textit{greater} efficiency in compute allocation compared to approaches relying on tokenization.

LLMs that utilize tokenization dedicate an equal amount of computational resources to each token. This approach compromises efficiency for performance because token boundaries are defined by compression strategies that do not always align with the prediction's difficulty. A key aspect of our model is the principle that computational effort should be assigned adaptively, according to necessity. For instance, the task of predicting the latter portions of most words generally involves straightforward, low-entropy choices, unlike predicting the initial word of a fresh sentence, which may require greater model capacity. This principle underpins the design of {\textsc BLT}'s architecture~(\S\ref{section:architecture}), which incorporates three transformer blocks: two compact, byte-level \textit{local models} alongside a single, large-scale global \textit{latent transformer}~(\autoref{fig:arch_high_level}). 

In order to flexibly assign computation and segment bytes into appropriate groups or patches, {\textsc BLT} bases segmentation on the entropy associated with predicting the following byte. This method yields contextually informed byte groupings, each exhibiting a comparable information density.

We introduce the first systematic flop-controlled scaling analysis of byte-level models extending to 8B parameters and 4T training bytes, demonstrating the feasibility of end-to-end large-scale training directly from bytes, eliminating the need for fixed-vocabulary tokenization. In terms of training efficiency measured by flops,\footnote{The computational requirement for each model is assessed by tallying the total number of Floating Point OPerations (flops) used.} {\textsc BLT} achieves parity with Llama 3 but requires up to 50\% fewer inference flops~(\S\ref{section:scaling}). Additionally, operating directly on raw bytes yields notable enhancements in modeling rare or less frequent patterns within the dataset. {\textsc BLT} models surpass their tokenizer-based counterparts in robustness to noisy input data and display superior character-level processing capabilities, as evidenced by their performance on tasks involving orthographic analysis, phonology, and low-resource machine translation~(\S\ref{sec:robustness}). Furthermore, {\textsc BLT} allows for simultaneous increases in both the model and patch size without exceeding a fixed inference flop allocation. Employing longer patch sizes generally results in computational savings, which can then be leveraged to expand the capacity of the global latent transformer, since it is executed fewer times. Through a series of scaling experiments constrained by a constant inference flop budget~(\autoref{fig:fixed_inference_scaling}), we find that byte-level models exhibit superior scaling properties when compared to architectures based on tokenization.

In conclusion, the contributions of this paper are as follows: 1) We present {\textsc BLT}, a byte-latent LLM framework that enables adaptive compute allocation to enhance flop efficiency, 2) We provide evidence that our approach achieves parity with Llama 3 on training flops up to the 8B parameter scale, while also allowing for up to 50\% improvements in flop efficiency with only marginal decreases in evaluation metrics, 3) The {\textsc BLT} architecture opens up a novel axis for scaling large language models, making it feasible to increase model size without exceeding a predefined inference cost, and 4) We empirically show that {\textsc BLT} models are more robust to noisy inputs and possess a deeper sensitivity to sub-word structures in data, aspects that are often overlooked by token-based LLMs. All code for training and inference of {\textsc BLT} is publicly available at~\url{https://github.com/facebookresearch/blt}.

\section{Patching: From Individual Bytes to Groups of Bytes}
\label{section:patching}

Dividing bytes into distinct \textit{patches} enables {\textsc BLT} to adjust its computational resources flexibly in response to the data’s context. Figure~\ref{fig:patching_types} illustrates multiple strategies for partitioning bytes into patches. More precisely, a patching function $f_p$ partitions an input byte sequence $\pmb{x}=\{x_i,|i=1,\ldots n\}$ of length $n$ into a series of $m < n$ patches $\pmb{p}=\{p_j|j=1,\ldots,m\}$ by assigning each $x_i$ a value in \{0,1\}, where a value of 1 signifies the beginning of a new patch. In both token-based and patch-based architectures, the main determinant of the required computational resources is the number of operations carried out by the core Transformer—specifically, the number of patches generated from a given patching function in {\textsc BLT}. Therefore, the mean number of bytes grouped within each patch, referred to as the \textit{patch size}, largely dictates the data processing cost at both training and inference time when using any specified patching method~(\S\ref{section:flops}).

Subsequently, we present three forms of patching functions: a fixed-length patching approach~(\S\ref{section:static-patch}), whitespace-delimited patching~(\S\ref{section:white-patch}), and adaptive patching based on entropy values obtained from a compact byte-level language model~(\S\ref{section:dyn-patch}). In closing, we outline approaches for incremental patching, as well as clarify distinctions between tokenization and patching~(\S\ref{section:bpe-patch}).

\subsection{Strided Patching Every K Bytes}
\label{section:static-patch}

A simple approach for organizing bytes is to divide them into uniform patches of size $k$, as implemented in MegaByte~\citep{yu2023megabyte}. This method's fixed stride offers convenience during both training and inference phases, allows easy adjustment of the mean patch length, and thus facilitates direct management of computational cost in terms of flops. Despite these advantages, this patching strategy has notable limitations. Firstly, computation is allocated uniformly, regardless of the content: as a result, a transformer step $j$ might be used inefficiently for predicting uninformative regions (e.g., whitespace in code), or conversely, may provide insufficient resources in segments with high information content, such as those containing mathematics. Secondly, such rigid partitioning produces inconsistent, non-adaptive patches for similar byte sequences—for example, instances of the same word might be divided differently across the data.

\subsection{Space Patching}
\label{section:white-patch}

According to \citet{slagle2024spacebyte}, a straightforward but effective advancement over traditional strided patching involves generating new patches immediately after encountering space-like bytes\footnote{
    Space-like bytes are characterized as any byte that does not correspond to a Latin letter, numeral, or UTF-8 continuation byte. Additionally, it is required that every patch includes at least one byte that is not space-like. }, which frequently align with linguistic unit boundaries in various languages. In the Space patching approach, each word triggers a latent transformer computation (requiring additional flops), thus dedicating modeling resources uniformly across all words within sequences. This approach ensures that words are consistently patched, and additional computation is directed to challenging predictions, which commonly occur at positions following spaces. For instance, determining the first byte in the answer to the question ``Who composed the Magic Flute?\uline{\hspace{.6em}}'' presents a much more complex task than predicting subsequent bytes after ``M'', since selecting the initial character substantially limits the set of viable completions, with the rest of ``Mozart'' being considerably more straightforward to infer. Nonetheless, space patching does not provide optimal handling across all languages and application areas, particularly because it lacks flexibility in patch size selection. In the subsequent section, we present an alternative patching approach, leveraging the understanding that word-initial bytes are typically the hardest to predict, while also offering an explicit means of adjusting patch size.

\subsection{Entropy Patching: Using Next-Byte Entropies from a Small Byte LM}
\label{section:dyn-patch}

Instead of utilizing a rule-based heuristic like whitespace, we employ a data-driven methodology to pinpoint next-byte predictions exhibiting elevated uncertainty. To this end, we propose \textit{entropy patching}, a technique that establishes patch boundaries based on entropy calculations.

We train a small byte-level auto-regressive language model on the training data for {\textsc BLT} and compute next byte entropies under the LM distribution $p_{e}$ over the byte vocabulary $\mathcal{V}$:

\begin{equation}
H(x_i) = \sum_{v \in \mathcal{V}} p_{e}(x_i=v|\pmb{x}_{<i}) \log p_{e}(x_i=v|\pmb{x}_{<i})
\end{equation}

We explore two approaches for detecting patch boundaries using the entropies $H(x_i)$. The first method selects points where the entropy surpasses a universal threshold, as depicted in~\autoref{fig:patching}. The second method seeks points whose entropy values are notably higher than the preceding value. This latter strategy may also be viewed as flagging points that disrupt an otherwise nearly monotonic decline in entropy within the patch.

\begin{align*}
    \text{Global Constraint}\hspace{1cm}H(x_t)& > \theta_g\\
    \text{Approx. Monotonic Constraint}\hspace{1cm}H(x_t)& - H(x_{t-1}) > \theta_r
\end{align*}

Patch boundaries are detected through a simple preprocessing procedure carried out in the data loading phase. This approach contrasts with the method proposed by~\citet{nawrot-etal-2023-efficient}, who employ a classifier trained to identify patch boundaries using entropy-based criteria. In our experiments~(\S\ref{section:experiments}), we perform a comparison of these two strategies for differentiating low-entropy from high-entropy bytes.

\subsection{The Byte-Pair Encoding (BPE) Tokenizer and Incremental Patching}
\label{section:incremental-patching}
\label{section:bpe-patch}

A large number of contemporary llms, such as our baseline Llama 3, rely on subword tokenizers like bpe~\citep{gage1994new, sennrich-etal-2016-neural}. In this context, we define ``tokens'' as byte sequences selected from a \textit{finite} vocabulary that is established before the start of training. This contrasts with ``patches,'' which are constructed by grouping sequences dynamically and do not depend on any fixed vocabulary. The essential distinction here is that, for tokens, the model does not have explicit access to the underlying byte-level representations.

A significant advancement of {\textsc BLT} compared to models that rely on tokenization is its reformulation of the balance between vocabulary magnitude and computational requirements. In conventional large language models, expanding the vocabulary results in longer tokens on average, thereby reducing the total number of model steps, but at the expense of increasing the dimensionality of the output in the model's final projection layer. This balance severely restricts tokenization-centric strategies in their ability to widely vary both token length and inference computational expense. As an illustration, Llama 3 increases the average number of bytes per token from 3.7 to 4.4, but this improvement is achieved by quadrupling the size of the embedding matrix compared to Llama 2.

During the generation process, {\textsc BLT} must determine at each step in the byte sequence whether it is currently at a patch boundary. This assessment is crucial because it dictates if additional computation by the Latent Transformer should be activated. Critically, this decision must be made without reliance on any bytes in the sequence that have not yet been generated; patching must proceed based only on existing information and cannot utilize knowledge of upcoming bytes to guide segmentation choices. More formally, a patching function $f_p$ is said to exhibit the property of incremental patching if the following holds:
$$f_p(\pmb{x}_{<i}) = f_p(\pmb{x})_{<i}$$
The bpe method does not fulfill the requirements of an incremental patching scheme, as the way a given prefix is tokenized by bpe can depend on its subsequent continuation; thus, bpe does not satisfy the above condition\footnote{One can enforce incremental patching in bpe by introducing a dedicated delimiter token to indicate patch boundaries, although this results in a longer byte sequence.}.

\section{{\textsc BLT} Architecture}
\label{section:architecture}
The {\textsc BLT} system consists of a large-scale global autoregressive language model, which processes patch-level representations, as well as two compact local models responsible for converting byte sequences into patches and for transforming patch representations back to bytes~(\autoref{fig:arch_high_level}).

\subsection{Latent Global Transformer Model}

\textit{The Latent Global Transformer} refers to an autoregressive transformer architecture, denoted as $\mathcal{G}$, comprised of $l_{\mathcal{G}}$ layers. This model processes sequences of latent input patch embeddings $p_j$ and generates corresponding sequences of output patch embeddings $o_j$. Throughout this work, the index $j$ is consistently used to reference patches, while $i$ indicates bytes. The global transformer employs a block-causal attention mask~\citep{dubey2024llama}, enforcing that each position attends only to its own and prior patches within a document. As the main computational component during both pre-training and inference, the usage of this model substantially determines total computational costs. Therefore, selecting when to utilize the global transformer enables flexible allocation of compute resources, allowing the system to adapt computational expenditure according to the complexity of different input and output segments.

\subsection{Local Encoder}
\label{section:local}

\textit{The Local Encoder Model}, symbolized by $\mathcal{E}$, employs a compact transformer architecture with significantly fewer layers ($l_{\mathcal{E}} << l_{\mathcal{G}}$) than its global counterpart. Its core purpose is to rapidly convert sequences of input bytes, $b_i$, into informative patch representations, $p_j$. Unlike standard transformers, this model introduces a cross-attention mechanism after each transformer block, which is responsible for aggregating the byte-level embeddings into higher-level patch representations~(\autoref{fig:crossattn}).

Initially, each input byte $b_i$ is mapped via an embedding matrix of size $\mathbb{R}^{256 \times h_{\mathcal{E}}}$, producing embedded vectors $x_i$. Optionally, these base embeddings can be supplemented with extra features through the integration of hash-embeddings~(\S\ref{sec:hash-emb}).

Subsequent to this, a stack of interleaved transformer and cross-attention layers~(\S\ref{sec:enc-cross-attn}) iteratively refines these representations, culminating in patch representations $p_i$ suitable for downstream processing by the global transformer, $\mathcal{G}$.

Each transformer block within $\mathcal{E}$ utilizes a \textit{local block causal} attention pattern, wherein each byte may attend only to a fixed span of $w_{\mathcal{E}}$ immediately preceding bytes; this window may span across dynamically determined patch divisions, but it is not permitted to extend beyond the bounds of the document. Details regarding the embedding procedure and the cross-attention mechanism are provided in the following subsections.

\subsubsection{Encoder Hash n-gram Embeddings}
\label{sec:ngram-emb}
\label{sec:hash-emb}
An essential aspect of producing rich and stable representations at every step $i$ is integrating contextual details from earlier bytes. In {\textsc BLT}, this is accomplished by capturing both the identity of the current byte $b_i$ in isolation \textit{as well as} its membership within a byte n-gram sequence. Specifically, at each position $i$, we generate byte-grams

\begin{align}
    g_{i,n}=\{b_{i-n + 1},\ldots, b_i\}
\end{align}

for each byte position $i$ and for each $n$ in the range from three to eight.\footnote{
We do not consider $n$-grams of length $n$ or longer if $i<n$. }

Subsequently, we define hash $n$-gram embeddings, wherein every byte $n$-gram is associated to a position within a dedicated embedding matrix $E_{n}^{hash}$, using a hash function, and with each $n \in\{3,4,5,6,7,8\}$ as discussed in \citet{bai2010learning}. The embedding obtained from this process is combined additively with the base embedding of the byte at position $i$; the result undergoes normalization prior to being delivered as input to the local encoder component. The composite embedding is computed by

\begin{align}
e_i &= x_i + \sum_{n = 3,...,8}E_{n}^{hash}(\text{Hash}(g_{i,n})) \\
\text{where, Hash}(g_{i,n}) &= \text{RollPolyHash}(g_{i,n}) \% |E_{n}^{hash}|
\end{align}

We scale $e_i$ by dividing by the total number of $n$-gram sizes incremented by one, and employ the RollPolyHash function, as specified in Appendix~\ref{appendix:rollpolyhash}. Section~\ref{section:ablations} presents an ablation study on how varying $n$-gram hash embeddings, specifically changing $n$ and the size of the embedding table, influences the trends in flop-controlled scaling laws. Alongside hash-based $n$-gram embeddings, we also conducted experiments utilizing frequency-based $n$-gram embeddings; comprehensive information regarding these trials is included in Appendix~\ref{appendix:ngrams}.

\subsubsection{Encoder Multi-Headed Cross-Attention}
\label{sec:enc-cross-attn}
Our approach largely mirrors the input cross-attention mechanism from the Perceiver architecture~\citep{jaegle2021perceiver}, but with a notable distinction: instead of utilizing a static set of latent representations, we align the latent variables with dynamic patch representations (\autoref{fig:crossattn}), restricting their attention exclusively to the bytes comprising each patch. Each patch $p_j$ has an associated query vector, initialized through a pooling operation applied to its corresponding byte embeddings, and then linearly projected by $\mathcal{E}_{C} \in \mathbb{R}^{h_{\mathcal{E}} \times (h_{\mathcal{E}} \times U_{\mathcal{E}})}$, where $U_{\mathcal{E}}$ represents the quantity of encoder cross-attention heads. Specifically, denoting the bytes of patch $p_j$ as $f_{\text{bytes}}(p_j)$, the computation proceeds as follows:

\begin{align}
P_{0,j} &= \mathcal{E}_{C}(f_\text{bytes}((p_j)), f~ \text{is a pooling function} \\
P_l &= P_{l-1} + W_o\left(\text{softmax}\left(\frac{QK^T}{\sqrt{d_k}}\right)V\right) \\
\text{where } Q_j &= W_q(P_{l-1,j}), K_i = W_k(h_{l-1,i}), V_i = W_v(h_{l-1,i}) \\
h_l &= \text{Encoder-Transformer-Layer}_l(h_{l-1})
\end{align}

Here, $P \in \mathbb{R}^{n_p \times h_{\mathcal{G}}}$ denotes the $n_p$ patch representations designated for input to the global model; these representations are initialized by aggregating the byte-level embeddings $e_i$ associated with each patch $p_j$. The matrices $W_q$, $W_k$, $W_v$, and $W_o$ are linear transformations for queries, keys, values, and output, where both keys and values are derived from projections of the byte representations $h_i$ from the preceding layer (and from $e_i$ in the initial layer). A patch-specific masking approach is implemented so that a query $Q_j$ only considers the keys and values related to the bytes within patch $j$. Given the use of multi-head attention across $Q$, $K$, and $V$, and recognizing that the patch representation typically has a larger dimension ($h_{\mathcal{G}}$) than $h_{\mathcal{E}}$, $P_l$ is preserved as several attention heads, each of size $h_{\mathcal{E}}$ during cross-attention, after which these representations are concatenated into vectors of dimension $h_{\mathcal{G}}$. Furthermore, we apply a pre-LayerNorm operation to queries, keys, and values, and omit positional embeddings in this cross-attention structure. A residual connection is also placed around the cross-attention module.

\subsection{Local Decoder} 

Analogous to the design of the local encoder, the local decoder $\mathcal{D}$ is implemented as a compact transformer architecture with $l_{\mathcal{D}}$ layers, where $l_{\mathcal{D}}$ is considerably smaller than $l_{\mathcal{G}}$. The decoder’s objective is to reconstruct a sequence of raw bytes, $y_i$, given a sequence of global patch representations $o_j$. Functioning autoregressively, the local decoder models each byte conditioned on the previously generated bytes. To accomplish this, it leverages the hidden states produced by the local encoder corresponding to the byte sequence. The processing involves $l_{\mathcal{D}}$ layers that alternate between cross-attention mechanisms and standard transformer operations. In each iteration, the cross-attention block precedes the transformer layer, ensuring that patch representations are initially mapped to byte-level representations, which are then refined by subsequent transformer layers within the local decoder.

\subsubsection{Decoder Multi-headed Cross-Attention} 

In the decoder cross-attention, the roles of the queries and key/values are interchanged i.e. the byte-representations are now the queries, and the patch representations are now the key/values. The initial byte-representations for the cross-attention are initialized as the byte embeddings from the last encoder layer i.e. $h_{l_{\mathcal{E}}}$. The subsequent byte-representations for layer $l$, $d_{l,i}$ are computed as:

\begin{align}
D_0 &= h_{l_{\mathcal{E}}} \\
B_l &= D_{l-1} + W_o\left(\text{softmax}\left(\frac{QK^T}{\sqrt{d_k}}\right)V\right), \\
  &\text{where } Q_i = W_q(d_{l-1,i}), K_i = W_k(\mathcal{D}_C(o_j)), V_i = W_v(\mathcal{D}_C(o_j)) \\
  D_l &= \text{Decoder-Transformer-layer}_l(B_l)
\end{align}

Here, $W_k$ and $W_v$ are matrices used for projecting keys and values, respectively; these projections are computed via a linear transformation, followed by a split operation $\mathcal{D}_C$ that processes the global model's terminal patch embeddings $o_j$. The matrix $W_q$ projects the queries, acting upon the byte-level representations $d_{l-1}$ produced by the previous decoder transformer layer (or $h_{l_{\mathcal{E}}}$, when computing the first decoder layer). $W_o$ serves as the matrix for the final output projection, so that $B$ has the shape $\mathbb{R}^{h_{\mathcal{D}} \times n_b}$, with $n_b$ being the total number of output bytes. The representations for the subsequent decoder layer, $D_l$, are produced by feeding the result $B$ from the cross-attention mechanism into a decoder transformer layer. Consistent with the approach in local encoder cross-attention, our configuration utilizes multi-head attention, pre LayerNorm, omits positional encodings, and includes a residual path encircling the cross-attention block.

\section{Experimental Setup}
\label{section:experiments}
We construct rigorously controlled experiments aimed at evaluating {\textsc BLT} alongside models employing traditional tokenization, ensuring that {\textsc BLT} does not benefit from potentially longer context windows in the comparison.

\subsection{Pre-training Datasets}  
In our experiments, models of all examined scales undergo pre-training using two main datasets: 1) The Llama 2 dataset~\citep{touvron2023llama}, which contains 2 trillion tokens sourced from a diverse array of publicly accessible materials, subsequently subjected to rigorous cleaning and filtering procedures to enhance data quality; and 2) {\textsc BLT}-1T: A newly assembled dataset consisting of 1 trillion tokens derived from assorted public repositories, which also incorporates a segment of pre-training content from Datacomp-LM~\citep{li2024datacomp}. The Llama 2 dataset serves as the basis for scaling law analyses—including token count optimization as investigated by~\cite{dubey2024llama}—to guide the selection of suitable architectural configurations for {\textsc BLT}. Conversely, {\textsc BLT}-1T is employed for a comprehensive pre-training cycle, enabling performance comparisons with Llama 3 across various downstream evaluation tasks. Both datasets strictly exclude any information sourced from Meta platforms or services. For baseline evaluations involving tokenizer-based models, we adopt the Llama 3 tokenizer—which comprises a 128K token vocabulary—as it yielded stronger performance baselines compared to the Llama 2 tokenizer in our assessments.

\subsection{Entropy Model} The entropy model employed in our study is a byte-level language model, trained using the same data distribution as the primary {\textsc BLT} model. By default, this model is implemented as a transformer with 100M parameters, comprising 14 layers and a hidden size of 512, and operates with a sliding window attention span of 512 bytes. All other hyperparameters match those used in both the local and global transformer variants. As outlined in \autoref{section:ablations}, we carried out experiments with various model capacities, receptive fields, and architectural choices. Notably, for models with sufficiently restricted receptive fields, the trained entropy model may be effectively represented as a compact lookup table.

\subsection{Entropy Threshold and Equalizing Context Length} For models employing entropy-based patching, we select a patching threshold that yields a targeted average \textit{patch size} across the pretraining dataset mix. In the case of {\textsc BLT}, the \textit{patch size} can be flexibly specified, as opposed to being determined by tokenization; this flexibility substantially influences the context length accessible to the model. To ensure fairness and prevent larger patches from providing undue benefit through longer contexts, we control for the expected number of bytes per batch, keeping it consistent across conditions. Consequently, when models utilize larger patch sizes, we proportionally decrease their sequence lengths. Specifically, with Llama 2 data, a context size of 8k bytes is used, whereas for the {\textsc BLT}-1T dataset, the context length is increased to an average of 16k bytes, all while maintaining an average batch size of 16M bytes.

Although the average batch size remains unchanged, dynamic patching approaches result in varying bytes-to-patch ratios during batch construction. To enhance efficiency, our implementation of {\textsc BLT} training assembles batches based on patches rather than bytes, thereby eliminating the need for padding within the more computationally intensive latent transformer. This batching strategy guarantees a uniform number of patches per batch. To manage memory usage and prevent spikes from exceptionally large patches, we pad or, if necessary, truncate input byte sequences to fixed lengths of 12k bytes for Llama 2 and 24k bytes for the {\textsc BLT}-1T datasets during training.

\subsection{Entropy Model Context}
\label{section:entropy-context}
Our observations indicate that entropy patching leads to increasingly larger patches in highly structured content, such as multiple choice assessments (as illustrated by the patching process on an MMLU instance in \autoref{fig:mmlu_patching}), which tend to be quite repetitive. These differences stem from reduced entropy present in the recurring elements within the entropy model’s context. Therefore, in the extensive implementation of {\textsc BLT}-Entropy with a patch size of 4.5, we refresh the entropy context by inserting new lines and employ an approximate monotonicity constraint, since this approach exhibits lower susceptibility to "entropy drift" associated with variations in context length. While this modification alters our method for calculating entropies, the approach for determining the entropy threshold remains unchanged.

\subsection{FLOPs Estimation} 
\label{section:flops}

Our approach to calculating transformer flops principally adheres to the methodology outlined in Chinchilla~\citep{hoffmann2022training}, which accounts for flops associated with the feed-forward network, the qkvo projections within the self-attention mechanism, as well as the computation of the attention and output projection steps. One important distinction in our procedure is the assumption that the input embedding layer utilizes an efficient lookup operation rather than a dense matrix multiplication, thereby rendering it a 0-flop process. Consistent with prior studies, we approximate that the backward pass consumes twice the computational resources (in flops) relative to the forward pass.

To compute flops \textit{per byte} for {\textsc BLT} models, we add up the flops for the local encoder transformer, the global latent transformer, and the local decoder transformer, together with the cross attention blocks in the encoder and the decoder:

\begin{align}
    \text{FL}_{\text{{\textsc BLT}}} &= \frac{1}{n_p}\text{Transf. FL}(h_{\mathcal{G}}, l_{\mathcal{G}}, m=n_{ctx}/n_p,V=0) \\
    &\quad + \text{Transf. FL}(h_{\mathcal{E}}, l_{\mathcal{E}}, m=w_{\mathcal{E}}, V=0) \\
    &\quad + \text{Transf. FL}(h_{\mathcal{D}}, l_{\mathcal{D}}, m=w_{\mathcal{D}}, V=256) \\
    &\quad + k/n_p \cdot \text{Cross Attn. FL}(h_{\mathcal{E}}, l_{\mathcal{E}}, m=n_p, r=n_p/k) \\
    &\quad + \text{Cross Attn. FL}(h_{\mathcal{D}}, l_{\mathcal{D}}, m=k, r=k/n_p)
\end{align}

Here, $n_{ctx}$ represents the input sequence length measured in bytes, while $n_p$ denotes the patch size. The parameter $r$ specifies the proportion between query and key/value representations. The value $k$ stands for the ratio of the patch dimension to the byte dimension and reflects how many local model splits are combined to produce a unified global model representation (with $k=2$ illustrated in \autoref{fig:crossattn}). The symbol $V$ signifies the size of the output vocabulary, which appears solely in the local decoder’s output projection layer. The variable $m$ refers to the context length and varies depending on whether attention is executed over a sequence of bytes or patches. Each component’s attention FLOPs are consequently adjusted to reflect its respective context length. Detailed formulations for both Transformer-FLOPs and Cross-Attention FLOPs calculations can be found in Appendix~\ref{section:flops-appendix}.

\subsection{Bits-Per-Byte Estimation} Perplexity only makes sense in the context of a fixed tokenizer as it is a measure of the uncertainty for each token. When comparing byte and token-level models, following previous work~\citep{xue2022byt5, yu2023megabyte, wang2024mambabyte}, we instead report Bits-Per-Byte (BPB), a tokenizer independent version of perplexity. Specifically:

\begin{align}
    \text{BPB}(x)&= \frac{\mathcal{L}_{CE}(\pmb{x})}{\ln(2)\cdot n_{\text{bytes}}}
\end{align}

Here, the total cross-entropy loss computed over the data $\pmb{x}$ is divided by the number of bytes in $\pmb{x}$ and a normalization factor, providing a normalized measure of the data’s uncertainty.

\subsection{Transformer Architecture Hyperparameters} 

For each transformer block within {\textsc BLT}, encompassing both local and global model components, we closely adopt the Llama~3~\citep{dubey2024llama} design. Specifically, the feed-forward networks incorporate the SwiGLU activation function~\citep{swiglu}, while self-attention modules utilize rotary positional embeddings (RoPE)~\citep{RoPe} set with $\theta=500000$ as recommended by~\citep{xiong2024effective}. Layer normalization is performed using RMSNorm~\citep{RMSNorm}. All self-attention layers that employ standard, static attention masks—such as \textit{block causal} and \textit{fixed-window block causal}—are implemented using Flash Attention~\citep{dao2022flashattention}, with a fixed attention window of size 512 for fixed-width masks. For cross-attention layers, where masking is dynamically determined based on patch dependencies, Flex Attention\footnote{\url{https://pytorch.org/blog/flexattention}} is utilized, enabling efficient fused kernel implementations that accelerate model training.

\subsection{{\textsc BLT}-Specific Hyperparameters} 

To evaluate the performance of {\textsc BLT} models, we design experiments focused on two main areas: scaling behavior and assessments on downstream tasks. Across these investigations, models of varying sizes are included: 400M, 1B, 2B, 4B, and 8B parameters. The detailed architectural settings for these configurations can be found in Appendix~Table~\ref{tab:arch}.

For initialization in the local encoder’s first cross-attention layer, query vectors are formed using max-pooling. All {\textsc BLT} models employ $500,000$ hashes generated with a single hash function, spanning n-grams of lengths 3 through 8. Training across all model sizes adopts a uniform learning rate of $4e-4$. This rate was identified as optimal through hyperparameter sweeps—comparing values between $1e-3$ and $1e-4$—conducted at both the 400M and 1B scales, where results demonstrated that both token and {\textsc BLT} models performed best under the same learning rate.

For scaling investigations with Llama-2 datasets, batch sizes are chosen following recommendations in~\cite{dubey2024llama} or appropriately converted to equivalent byte-size. Optimization is carried out using the AdamW optimizer~\citep{adamw}, with $\beta_1$ set to 0.9 and $\beta_2$ set to 0.95; we set $\epsilon=10^{-8}$. Training incorporates a warm-up phase linearly ramped over 2000 iterations and follows with cosine annealing of the learning rate down to zero. Additionally, a weight decay parameter of 0.1 is applied, and the gradient norms are clipped globally at 1.0.

\section{Scaling Trends}
\label{section:scaling}

We provide a comprehensive overview of scaling behavior in byte-level models, with the objective of guiding future scaling of {\textsc BLT} models. Our investigation into scaling distinguishes itself from earlier studies on byte-level models in several important aspects: (a) we assess performance patterns within the compute-optimal training framework, (b) we develop comparable 8B-parameter models trained on substantial datasets—reaching as much as 1T tokens (equivalent to 4T bytes)—and subsequently benchmark these models on a variety of downstream tasks, and (c) we evaluate scaling properties under scenarios where inference cost is held constant. In subsequent sections, we delve deeper into unique benefits that arise from the modeling of byte sequences.

\subsection{Parameter Matched Compute Optimal Scaling Trends}
\label{section:parameter-matched-scaling}
Employing the Llama 2 dataset, we train several \textit{compute-optimal} bpe and {\textsc BLT} models at four scales, spanning from 1B up to 8B parameters. We visualize the relationship between the training flops and language modeling accuracy for a representative sample of the overall training data mixture. The bpe models adhere to the model-to-data ratio found optimal in Llama~3~\citep{dubey2024llama}. This \textit{compute-optimal} methodology is theoretically formulated to yield maximal performance on the training data for a prescribed computational expenditure~\citep{hoffmann2022training}, establishing a reliable point of comparison for our analyses. In parallel to each bpe model, we also construct a {\textsc BLT} counterpart using the identical dataset, with each Latent Transformer configured to match the architecture and size of its bpe Transformer equivalent.

As shown in~\autoref{fig:scaling-compute-opt} (right), {\textsc BLT} models consistently equal or surpass their bpe equivalents, with this pattern persisting as both the model scale and computational requirements increase. To our knowledge, {\textsc BLT} represents the first instance of a byte-level Transformer architecture exhibiting similar scaling performance to BPE-based models within compute-efficient settings. This observation supports our hypothesis that the optimal balance between model parameters and training compute observed for bpe extends to {\textsc BLT}, or is at least very closely aligned.

Both advancements in architecture and the implementation of dynamic patching are essential for achieving the scaling characteristics observed with bpe models. As shown in ~\autoref{fig:scaling-compute-opt} (left), we benchmark space-patching-based models relative to Llama 3. Here, SpaceByte \citep{slagle2024spacebyte} is estimated using {\textsc BLT} space-patching, deliberately excluding n-gram embeddings and cross-attention. While SpaceByte represents a notable step up from Megabyte, its performance still significantly lags behind that of Llama 3. In \autoref{fig:scaling-compute-opt} (right), we depict how both structural refinements and dynamic patching jointly boost performance. At scale, {\textsc BLT} models are competitive with leading tokenizer-based architectures, including Llama 3.

Additionally, our results indicate that for tokenizer-based models, selecting the Llama-3 tokenizer yields superior performance compared to utilizing the Llama-2 tokenizer, even when both are trained on identical datasets.

Ultimately, our {\textsc BLT} architecture exhibits behavior that falls between that of Llama 2 and 3, particularly when much larger patch sizes are utilized. While the bpe tokenizers employed by Llama 2 and 3 yield an average token size of 3.7 and 4.4 bytes respectively, {\textsc BLT} demonstrates similar scalability with average patch sizes as large as 6 and even 8 bytes. Since inference flop count is inversely related to average patch size, opting for a patch size of 8 bytes translates to almost a 50\% reduction in inference flops. Additionally, models configured with greater patch sizes appear to benefit more as model capacity and data volume are increased. Specifically, although {\textsc BLT} with an 8-byte patch size performs noticeably worse than bpe Llama 2 at the 1B parameter scale, its performance surpasses bpe at the 7B scale. This pattern implies that adopting larger patch sizes could be even more advantageous as model and dataset sizes continue to expand, and that scaling up patch sizes further may be feasible as available compute and model scale increase.

\subsection{Beyond Compute Optimal Task Evaluations}
\label{section:evals}
To further investigate scaling behaviors, we continue training an 8B {\textsc BLT} model beyond the compute-optimal ratio, utilizing the more extensive and higher-quality {\textsc BLT}-1T dataset. We then evaluate its abilities using a broad set of standard benchmarks that test classification and generation skills. The chosen tasks emphasize common sense reasoning, general knowledge, and code synthesis:
\paragraph{Classification tasks} involve ARC-Easy (0-shot)~\citep{Eval_ARC}, Arc-Challenge (0-shot)~\citep{Eval_ARC}, HellaSwag (0-shot)~\citep{Eval_hellaswag}, PIQA (0-shot)~\citep{Eval_piqa}, and MMLU (5-shot)~\citep{hendrycks2020measuring}. For assessment, we use a prompt-based likelihood evaluation over the available answer options, computing and reporting mean accuracy scores.
\paragraph{Coding related generation tasks:}  To assess the model's Python code generation capabilities, we present pass@1 results on MBPP (3-shot)~\citep{Eval_MBPP} and HumanEval (0-shot)~\citep{Eval_HumanEval}.

In \autoref{tab:evals}, we present a comparative analysis of three models trained on the {\textsc BLT}-1T dataset: a model based on the Llama 3 tokenizer with bpe,\footnote{The Llama 3 tokenizer, featuring a 128k vocabulary, was selected due to its superior performance over the 32k vocabulary of Llama 2.} along with two versions of the {\textsc BLT} model. The first variant applies a space-patching approach ({\textsc BLT}-Space), while the second employs an entropy-driven patching mechanism ({\textsc BLT}-Entropy), which imposes an approximate monotonicity constraint and resets the entropy model's context with new lines, as described in~\autoref{section:entropy-context}. Each model is trained under an identical computational (flop) budget. For the {\textsc BLT}-Entropy variant, we further lower the entropy threshold at inference from 0.6 to 0.1, an adjustment that enhances performance on various tasks but leads to a greater number of inference steps.

The {\textsc BLT}-Entropy model achieves superior results compared to the Llama 3 model in 4 of the 7 evaluated tasks, despite both being trained on an equivalent byte count. This advantage likely stems from two factors: (1) more efficient utilization of training resources through dynamic patching, and (2) explicit modeling at the byte level rather than at the token level.

Conversely, {\textsc BLT}-Space lags behind the Llama 3 tokenizer in performance across all tasks except one. Nevertheless, it offers notable savings in inference flops due to its increased average patch size of 6 bytes. By contrast, both the bpe model and those employing entropy-patching achieve similar average patch sizes of about 4.5 bytes on the combined training dataset. Given an equivalent training budget, the model with the larger patch size is able to process 30\% more data than the other two, which could drive {\textsc BLT} away from the compute-optimal regime.

\subsection{Patches Scale Better Than Tokens} In {\textsc BLT} models, both model capacity and patch dimensions can be scaled up together without exceeding the available training and inference FLOPs or altering the volume of training data. The distinctive capability of patch-based architectures to flexibly increase patch size allows them to overcome the efficiency limitations inherent to fixed-vocabulary, token-based approaches, as outlined in Section~\ref{section:bpe-patch}. Extending patch lengths reduces the overall computational cost, enabling more resources to be devoted to expanding the global latent transformer, since this component is invoked less frequently.

We perform a fixed inference scaling study to evaluate whether utilizing larger models that take fewer computational steps over bigger patches yields superior results compared to using smaller models with more steps. Using the Llama 2 tokenizer, we begin with model configurations of 400 million and 3.6 billion parameters and identify flop-equivalent models employing both the Llama 3 tokenizer and {\textsc BLT}-Entropy models with mean patch sizes of 6 and 8 bytes on the training datamix (refer to~\autoref{tab:fixed-inf-params} for comprehensive model specifications). In the case of patch size 8, three encoder layers are deployed rather than one. Each model undergoes training with several different training flop constraints.

\autoref{fig:fixed_inference_scaling} illustrates that {\textsc BLT} architectures exhibit more favorable scaling behavior compared to tokenization-based models across both types of inference flop classes. While BPE models outperform at lower training compute allocations, {\textsc BLT} surpasses them once training extends beyond the compute-optimal threshold. In real-world scenarios, allocating additional resources to pretraining can be advantageous, as it allows for superior model performance under a constrained inference compute budget. This is exemplified by 8B models such as Llama 3.1, which underwent training on data volumes two orders of magnitude greater than the compute-optimal amount for its parameter scale.

The threshold at which {\textsc BLT} surpasses token-based models has moved marginally nearer to the compute-optimal point as model size increases within the higher flop category (shifting from approximately 3x to 2.5x the compute-optimal allocation). Furthermore, the variant with a patch size of 8 in this higher flop class demonstrates a more pronounced scaling behavior, allowing it to outperform the other models at an earlier stage. Referencing Section~\ref{section:parameter-matched-scaling}, it is evident that models with larger patch sizes achieve performance closer to BPE models as model scale expands. We partly attribute this effect to the reduced proportion of overall flops consumed by the byte-level Encoder and Decoder components, since their computational requirements seem to grow more slowly compared to the Latent Transformer. When the total parameter count is increased 20-fold—from 400M to 8B—{\textsc BLT}'s local model parameters only grow by roughly a factor of two. This distinction matters, since using larger patch sizes only impacts the computational cost for the patch Latent Transformer, not the byte-level parts. Indeed, this explains the rise of {\textsc BLT}-Entropy ps=8 from 1.6x to 1.7x the Llama 2 model size as the model scale increases.

To summarize, our investigation into patch-length scaling reveals that the {\textsc BLT} patch-based architecture benefits from increasing both patch and model size at the same time, yielding superior scaling patterns. These observed trends not only remain stable but appear to strengthen as model size grows.

\section{Byte Modeling Improves Robustness} We further evaluate the robustness of {\textsc BLT} relative to token-based counterparts that do not incorporate explicit byte-level data, and introduce a methodology for converting pretrained token-based models to operate at the byte level.
\label{sec:robustness}

\subsection{Character-Level Tasks}

One of the initial reasons for developing byte-level models was their enhanced resilience to noise at the byte level within input sequences, as well as their capability to recognize the internal composition of tokens—a known challenge for existing tokenizer-based approaches. To assess these aspects, we conduct supplementary experiments on datasets specifically designed to test both robustness against noisy inputs and the ability to recognize characters in English and other languages, encompassing numerals and phonemic elements. The outcomes of these analyses are summarized in Table~\ref{tab:char_tasks}.

\paragraph{Noisy Data} To assess and compare the robustness of {\textsc BLT} against tokenizer-based models, we construct noisy variants of the benchmark classification tasks outlined in Section~\ref{section:evals}. We implement five unique character-level noise techniques to generate diverse textual perturbations: (a) \textit{AntSpeak}, in which the complete text is rendered in uppercase and each character is separated by a space; (b) \textit{Drop}, which randomly eliminates 10\% of all characters; (c) \textit{RandomCase}, where half the characters are capitalized and the other half are converted to lowercase at random positions throughout the text; (d) \textit{Repeat}, whereby 20\% of characters are repeated up to four consecutive times; (e) \textit{UpperCase}, which changes the entire text to uppercase. During evaluation, we independently apply each noise type to the prompt, completion, or both, and subsequently average the resulting scores. Table \ref{tab:char_tasks} shows performance metrics for noised HellaSwag~\citep{Eval_hellaswag}, demonstrating that {\textsc BLT} consistently surpasses tokenizer-based models in terms of noise robustness, achieving an average improvement of 8 points over a similarly-trained model, and even outperforming the Llama 3.1 model that had access to substantially more training data.

\paragraph{Phonology - Grapheme-to-Phoneme (G2P)} We evaluate how effectively {\textsc BLT} translates a sequence of graphemes—written symbols corresponding to a word—into its phonemic representation, which encodes pronunciation. Table \ref{tab:char_tasks} displays the outcomes of the G2P experiment conducted in a 5-shot regime, utilizing Phonology Bench~\citep{suvarna2024phonologybench}. The results indicate that {\textsc BLT} achieves higher performance than the Llama 3 1T tokenizer-based baseline model for this task.

\paragraph{CUTE}
To measure the ability of models to process text at the character level, we evaluate {\textsc BLT} on the CUTE benchmark~\citep{edman2024cute}. This benchmark encompasses various tasks divided into three main groups: tasks evaluating compositional understanding, those testing orthographic similarity, and tasks focused on sequence manipulation. The CUTE benchmark presents significant difficulties for most models relying on tokenization, as such models typically encode information about token spellings but often fail to leverage it for text manipulation. As shown in Table~\ref{tab:char_tasks}, {\textsc BLT}-Entropy surpasses both BPE Llama 3 models by over 25 points on this benchmark. Notably, the model shows near-perfect performance on tasks involving manipulation at the character level, reaching 99.9\% accuracy in both spelling-related evaluations. These substantial gains, achieved even though {\textsc BLT} was trained with 16 times less data than Llama 3.1, suggest that BPE-based models have particular difficulty in learning character-level representations. Figure \ref{fig:cute_examples} displays examples where the Llama 3 tokenizer-based model falters, but our {\textsc BLT} model succeeds. The only exceptions are the word deletion and insertion tasks, where BPE-based approaches perform better; this might reflect inherent challenges byte-level models face in these specific operations. However, the performance gap in these areas is relatively minor, and building word-level understanding from characters may be less complex than extracting character-level representations from words. For all evaluations, we employ a consistent experimental setup and use the original prompts from Huggingface. It should be noted that BPE models might see improvement from further prompt refinement.

\paragraph{Low Resource Machine Translation} We assess the performance of {\textsc BLT} on translation tasks involving six widely used language families and twenty-one low-resource languages, each employing diverse scripts, using the FLORES-101 benchmark~\citep{goyal2022flores}. SentencePiece BLEU scores are presented in Table~\ref{tab:flores}.
The findings indicate that {\textsc BLT} surpasses a model utilizing the Llama 3 tokenizer, exhibiting an average improvement of 2 BLEU points for translations into English and a 0.5-point gain for translations out of English. For common language pairs, {\textsc BLT} achieves results on par with or marginally superior to those of Llama 3. Notably, {\textsc BLT} consistently outperforms Llama 3 across many pairs within low-resource language families, illustrating the advantage of byte-level modeling when addressing long-tail byte sequence generalization.

\subsection{Training {\textsc BLT} from Llama 3}
\label{sec:distillation}
We investigate an approach wherein {\textsc BLT} models benefit from pre-existing tokenizer-based models to accelerate and enhance training outcomes. This is accomplished by initializing the global transformer parameters of {\textsc BLT} using those extracted from a pre-trained Llama 3.1 model. Afterwards, the global transformer weights are fine-tuned with a learning rate set at one-tenth of that used for both the local encoder and decoder, matching the recommended training duration for Llama 3. Results comparing this setup to a baseline {\textsc BLT} are reported in Table~\ref{tab:distill}. These findings demonstrate that initializing {\textsc BLT} with Llama 3.1 parameters yields performance that not only surpasses both standard Llama 3 and {\textsc BLT} trained under equivalent computational budgets, but also exceeds the results of our {\textsc BLT}-Entropy model on the MMLU benchmark (see Table~\ref{tab:evals}), even though the latter was trained with a much larger corpus (1T tokens). This indicates the potential of this strategy for notable reductions in total training flops while maintaining or improving downstream effectiveness.

This approach can be interpreted as adapting tokenizer-based architectures into models that do not require a tokenizer, essentially turning a pre-trained LLaMA 3.1 system into a {\textsc BLT} model. For a thorough comparison, we present results for the original LLaMA 3.1 trained on 15T tokens in Table \ref{tab:distill}, evaluating it alongside the {\textsc BLT} obtained from LLaMA 3. Our results indicate a modest decrease in performance on MMLU and HumanEval, with more pronounced degradation on other benchmarks. These outcomes imply that there is room for advancement in fully exploiting pre-trained weights, particularly through refining dataset composition and tuning additional hyperparameters.

\section{Ablations and Discussion}
\label{section:ablations}
Here, we present ablation studies to motivate the architectural selections made for {\textsc BLT}, as well as to elaborate on the chosen patching strategy and the set of hyper-parameters used in training the 8B parameter {\textsc BLT} model on the {\textsc BLT}-1T dataset.

\paragraph{Entropy Model Hyper-parameters} In order to investigate how adjustments to entropy model size and the length of the context window influence scaling behavior, we conduct experiments using byte-level entropy transformer architectures. These models are configured across a range of parameter counts, from 1 million up to 100 million, and with context window sizes spanning from 64 to 512 tokens. We depict the relationship between bits-per-byte (bpb) and the amount of training computational resources (FLOPs) on scaling law plots, leveraging our $400m$ and $1b$ {\textsc BLT} models, which have been trained on the Llama-2 dataset; the resulting curves are shown in \autoref{fig:entropy_ablation}. Our analysis indicates that larger model sizes and longer context windows are generally beneficial for scaling performance, though increases beyond 50 million parameters yield progressively smaller improvements.

\paragraph{Types of Patching}

We conduct ablation studies on the four patching methods described in Section~\ref{section:patching}, namely: (1) Strided Patching utilizing strides of 4 and 6, (2) Whitespace-based Patching, (3) BPE Tokenizer patching utilizing the Llama 3 tokenizer, and (4) Entropy-driven patching employing a lightweight byte-level LLM.

Although dynamic patching decreases the actual sequence length, we regulate sequence length to ensure that every patching approach encounters a comparable context length. This way, during both training and inference, each model is exposed to an equivalent expected number of bytes per sequence, thereby mitigating the risk of confounding variables related to modeling extended contexts. The outcomes of these ablation studies are presented in Figure~\ref{fig:scaling-compute-opt}. In these results, all alternative patching methods surpass static patching in performance, and notably, space patching nearly matches the effectiveness of dynamic entropy-based patching.

In \autoref{tab:evals_ablation}, we show benchmark results for {\textsc BLT} models that contrast tokenizer-based models, space patching, and entropy-based patching approaches, all trained on the Llama 2 dataset for the optimal step count~\citep{dubey2024llama}. While space patching represents a more straightforward approach, as it avoids running an entropy model dynamically during training, our findings demonstrate that the improvements yielded by entropy-based patching—previously identified in scaling experiments~(Section~\ref{section:scaling})—persist across a variety of downstream benchmark tasks.\footnote{Space patching outcomes are from earlier experiments conducted without cross-attention, but consistent patterns emerge when cross-attention is incorporated.}

\paragraph{Cross-Attention}
In~\autoref{tab:crossattn_ablations_1b}, we examine the effects of incorporating cross-attention at different locations within both the encoder and decoder components of {\textsc BLT}. For cross-attention in the encoder, we explore three initialization strategies for the queries: 1) a universally learned embedding applied across all global states, 2) a hash embedding derived from the bytes within each patch, and 3) pooling the encoder's hidden representations of the patch bytes at a specific encoder layer.

Our results indicate that integrating cross-attention within the \textit{decoder} yields the strongest performance gains. Implementing cross-attention in the encoder leads to only modest improvements, which are evident primarily when the queries are initialized via pooling. Furthermore, we observe that the benefits of cross-attention are pronounced on the Common-Crawl dataset and become increasingly significant when employing larger patch sizes.

\paragraph{n-gram Hash Embeddings} 

We examine the impact of various n-gram hash embedding vocabulary sizes—specifically, 0, 100K, 200K, and 400K—and summarize our observations in Table~\ref{tab:ngram_ablations_1b}. Our analysis shows that incorporating hash embeddings consistently benefits all evaluated domains, with particularly pronounced improvements on Wikipedia and Github (yielding a 0.04 bpb improvement as opposed to a 0.01 bpb improvement after 15k steps at the 8B scale). When increasing hash size from 300K to 500K at the 8B model size, the resulting change in performance at 15k steps was minimal (approximately 0.001 bpb). These findings suggest that hash embeddings play a crucial role in narrowing the performance gap between {\textsc BLT} models and those utilizing traditional tokenizers. However, we also observe that returns tend to plateau after reaching 300K hashes. Furthermore, the performance enhancements from hash embeddings and cross-attention mechanisms appear largely additive, as they contribute to distinct improvements across multiple datasets.

\paragraph{Local Model Hyperparameters} Table \ref{tab:local_layers} presents an ablation study exploring different configurations for the number of layers in both the local encoder and decoder. Notably, when utilizing hash n-gram embeddings, {\textsc BLT} achieves strong performance with a highly minimal encoder—consisting of only a single layer—while employing a more substantial decoder.

\section{Related Work}
\label{section:related_work}

\textbf{Character-Level RNNs:} The task of Character Language Modeling has long held significance in the neural modeling landscape~\citep{sutskever2011generating,mikolov2012subword,graves2013generating}, largely due to its inherent ability to seamlessly handle out-of-vocabulary words without depending on explicit back-off strategies. In their work, \cite{kim2016character} developed a framework that leverages convolutional and highway network architectures on the input side—processing only characters—and subsequently connects these representations to LSTM-based RNNs. This architecture achieved parity with then state-of-the-art RNN language models for English, while surpassing them on languages characterized by rich morphological structures, thereby highlighting a key strength of character-level LLMs. Furthermore, \cite{kenter2018byte} introduced a byte-level LSTM-based machine comprehension approach, which demonstrated superior performance over word-based models, notably on morphologically complex languages such as Turkish and Russian. Complementarily, \cite{zhang2015character} explored the application of character-level convolutional models to classification problems, finding that for certain tasks, these models outperformed their word-based counterparts. In a similar vein, \citet{chung2019hierarchical} proposed hierarchical LSTM models that utilize boundary-detectors at various levels to autonomously infer the underlying hierarchical structure in text, thereby boosting results in character-level language modeling. Diverging from attention mechanisms, ByteNet~\cite{kalchbrenner2016neural} adopts convolutional layers directly over character sequences for the purpose of machine translation.

\textbf{Character-Level Transformers:} The advent of transformer architectures leveraging attention mechanisms~\citep{vaswani2017attention} in conjunction with subword tokenization techniques~\citep{sennrich-etal-2016-neural} brought substantial advancements to language modeling and standard NLP benchmarks. Nevertheless, employing word or subword segments inherently imposes an inductive bias regarding the linguistic abstraction at which models function. Seeking to integrate the advantages of transformers with the encouraging early outcomes of character-based language modeling, \citet{al2019character} experimented with extremely deep transformer networks, utilizing auxiliary loss functions to enable transformer-based systems that achieved superior performance compared to prior character-level LSTM models. Despite these gains, their models continued to underperform relative to word-level large language models. Similarly, GPT-2~\citep{radfordgpt2} noted that, on massive corpora such as the 1 Billion Word Benchmark, models operating at the byte level lagged behind those trained with word-level representations.

Although \cite{Choe2019BridgingTG} showed that transformer-based language models operating at the byte level can surpass subword-level counterparts given similar parameter counts, such architectures demand substantially greater computational resources and longer training times. Along similar lines, \cite{el2020characterbert} introduced the CharFormer, a BERT-style model where word representations are constructed by applying convolutional networks to character embeddings. Their approach yielded domain-specific improvements in medical applications, albeit at a considerable computational cost. The work by \cite{clark2022canine} presented CANINE, a 150M-parameter encoder-only model designed for character-level processing. Structurally, CANINE employs a deep transformer as its main component, akin to our global model, while leveraging a combination of local transformers and strided convolutions to efficiently compress sequences of input characters. This model outperformed the token-based encoder-only baseline (mBERT) on a variety of multilingual downstream tasks. ByT5~\citep{xue2022byt5} explored methods for constructing byte-level encoder-decoder models that eschew any patching mechanisms. Despite the increased robustness of ByT5 models to input noise and their competitiveness with tokenized models—achieving similar performance using four times less data—the absence of patching necessitated resource-intensive attention computations across all bytes, substantially raising the computational load. Shifting from subword to byte-level modeling inflates input sequence lengths, thereby complicating the task of scaling such models efficiently. More recently, leveraging the Mamba Architecture~\citep{gu2023mamba}, which is capable of retaining a fixed-size memory over very lengthy contexts, \cite{wang2024mambabyte} developed a byte-level model without patching and demonstrated superior performance in a FLOP-matched comparison with byte-level transformer models at the 350M parameter scale, achieving better bits-per-byte scores on multiple benchmarks.

\textbf{Patching-based approaches:} The strategy of employing patching methods has proven useful in significantly reducing the excessive computational overhead, measured in FLOPs, associated with byte-level LLMs while still potentially maintaining robust performance. Multiple studies have highlighted the initial effectiveness of these methods with smaller models and training data regimes. For example, \cite{nawrot-etal-2022-hierarchical} introduce a static patching approach involving downsampling and upsampling, culminating in the hourglass transformer architecture, which surpasses alternative byte-level baselines at the 150M parameter scale. Building upon this, \cite{nawrot-etal-2023-efficient} integrate dynamic patching strategies such as a learned, end-to-end boundary predictor, a boundary predictor trained with guidance from particular tokenizers, and an entropy-driven patching mechanism akin to {\textsc BLT}. Their findings indicate that these dynamic schemes yield superior performance compared to contemporaneous vanilla transformers for language modeling tasks on approximately 400M tokens with a model containing 40M parameters. In parallel, \cite{lester2024training} focus on training with sequences that are compressed via arithmetic coding, reaching levels of compression beyond those typically achieved with BPE. By adopting an equal-info windows methodology, they demonstrate superior results relative to byte-level baselines on language modeling benchmarks, though their approach does not match the results of subword-based models.

Our research is primarily motivated by MegaByte~\citep{yu2023megabyte}, a decoder-only causal LLM that implements fixed, static patching along with concatenation of representations to map bytes into patches, utilizing a local model on the decoding end to revert patches back to bytes. The MegaByte framework was shown to rival the performance of models employing tokenizers at the 1B parameter level, particularly on datasets comprising roughly 400B bytes. In our experiments, we systematically evaluate MegaByte and observe that, within settings where FLOPs are controlled, the static patching approach does not reach the performance of the latest compute-efficient, tokenizer-based models. We further show that {\textsc BLT} effectively narrows this disparity. In parallel, \cite{slagle2024spacebyte} echo similar conclusions regarding MegaByte and propose modifications—such as applying patching based on whitespaces and analogous byte-types, in addition to incorporating a local encoder component. They report gains over tokenizer-based transformers in certain contexts like Github and arXiv for models with 1B parameters. We also present empirical results for this variant, emphasizing that additional innovations in architecture are essential for byte-level models to scale further and to ultimately reach parity with cutting-edge token-based architectures such as Llama 3.

\section{Limitations and Future Work}

In this study, we adopt model training regimes based on the optimal number of steps established for Llama~3~\citep{dubey2024llama} when making architectural decisions. It should be noted, however, that these scaling relationships were originally derived for BPE-level transformers, which may result in non-ideal (data, parameter size) proportions when applied to {\textsc BLT}. Investigating scaling laws tailored specifically to {\textsc BLT} is a promising direction for future research, as these may reveal even more advantageous scaling patterns for our model. Furthermore, our experiments are primarily limited to the 1B parameter range; as model size increases to 8B parameters or more, optimal architecture configurations may also shift, potentially leading to enhanced performance at greater scales.

Current transformer libraries and codebases are primarily optimized for high efficiency with tokenizer-based transformer models. Although we conduct experiments with theoretical FLOP matching and incorporate some optimized implementations (like FlexAttention) to manage components that differ from the standard transformer, our models may still lag behind tokenizer-based variants regarding wall-clock performance and could achieve additional gains through further optimization efforts.

Whereas {\textsc BLT} employs an independently trained entropy model for the purpose of patching, an intriguing avenue for further research would involve integrating patching model learning into an end-to-end training pipeline. Section \ref{sec:distillation} describes preliminary results, which suggest that ``byte-ifying'' large tokenizer-based architectures such as Llama 3—models trained on datasets exceeding 10T tokens—is viable by initializing and freezing the global transformer with their pretrained parameters. Expanding upon these initial findings may lead to techniques that preserve the strengths of bytefying while also surpassing the performance of the original tokenizer-based systems, all without the necessity for full retraining.

\section{Conclusion}
\label{section:conclusion}

This work introduces the Byte Latent Transformer (\textbf{BLT}), an innovative model architecture that challenges the standard reliance on fixed-vocabulary tokenization prevalent in large language models. Through the implementation of a trainable, adaptive mechanism for aggregating bytes into variable-sized patches, {\textsc BLT} allocates computation in accordance with the underlying data’s complexity, resulting in marked gains in both efficiency and resilience. Comprehensive scaling analyses show that {\textsc BLT} can reach parity with token-based systems such as Llama 3 for models up to 8B parameters and 4T bytes, and can achieve up to 50\% lower inference flops at the cost of only slight decreases in standard benchmarks. Moreover, {\textsc BLT} enables concurrent scaling of both model size and patch size while keeping inference computational demands constant, which is particularly advantageous for the computational constraints typically found in real-world applications. By processing raw byte sequences directly, {\textsc BLT} further enhances the model’s robustness when faced with irregular or noisy data and provides greater insight into sub-word level structures. Collectively, these findings highlight {\textsc BLT} as a compelling and scalable alternative to conventional tokenization-based models, supporting more resource-efficient and adaptive language processing.

\end{document}